% Created 2026-05-11 Mon 12:30
% Intended LaTeX compiler: pdflatex
\documentclass[11pt]{article}

\usepackage[utf8]{inputenc}
\usepackage[T1]{fontenc}
\usepackage{graphicx}
\usepackage{longtable}
\usepackage{wrapfig}
\usepackage{rotating}
\usepackage[normalem]{ulem}
\usepackage{amsmath}
\usepackage{amssymb}
\usepackage{capt-of}
\usepackage{hyperref}
\usepackage{minted}
\author{flavio}
\date{\today}
\title{Eridetarieta}
\hypersetup{
 pdfauthor={flavio},
 pdftitle={Eridetarieta},
 pdfkeywords={},
 pdfsubject={},
 pdfcreator={},
 pdflang={English}}
\begin{document}

\maketitle
\tableofcontents

\section{Eriditarietà}
\label{sec:org405a5de}
L'eriditarietà è una forma di riuso del software in cui una classe è creata e in cui assorbiamo i membri di una classe esistente. Possiamo anche aggiungere nuove caratteristiche o migliorare quelle esistenti. Diciamo che una classe eredità da un'altra usando la parola chiave \texttt{extends}
\subsection{Esempio}
\label{sec:org002d543}
\begin{minted}[]{java}
public class Forma
    {

        public void disegna(){}
}
\end{minted}

\begin{minted}[]{java}
public class Triangolo extends Forma{
    private double base;
    private double altezza;
    public Triangolo(double base, double altezza){
        this.base = base;
        this.altezza = altezza;
    }
    public double getBase(){
        return base;
    }
    public double getAltezza()
        {
           return altezza;
        }

public static void main(String[] args){
    Triangolo t = null;
}
    }
\end{minted}
\end{document}
